\lyseis{5.2}{Προβλήματα με ποσοστά}
\sol{1} α) $0{,}8\cdot45=36$ € \quad β) $9{,}6$ € \quad γ) $104$ €.
\sol{2} α) $1{,}15\cdot40=46$ € \quad β) $230$ € \quad γ) $20{,}7$ €.
\sol{3} α) $\frac{10}{50}=20\%$ αύξηση \quad β) $\frac{20}{80}=25\%$ μείωση \quad γ) $\frac{10}{25}=40\%$ αύξηση \quad δ) $\frac{30}{120}=25\%$ μείωση.
\sol{4} α) $1{,}2x=72$, $x=60$ \quad β) $0{,}6x=45$, $x=75$ \quad γ) $0{,}15x=12$, $x=80$.
\sol{5} Για $248$ €: $\frac{24}{124}\cdot248=48$ € ΦΠΑ και αξία $200$ €. Για $62$ €: ΦΠΑ $12$ € και αξία $50$ €.
\sol{6} α) $0{,}04\cdot1\,500=60$ € \quad β) $30$ €.
\sol{7} Με $30\%$ και $20\%$: $0{,}7\cdot0{,}8\cdot90=50{,}40$ €. Με $50\%$: $45$ €. Διαφορά $5{,}40$ €.
\sol{8} Μπουφάν: $0{,}85\cdot120=102$ €. Μπότα: $0{,}8\cdot85=68$ €. Μαζί $170$ €.
\sol{9} $250\cdot1{,}1=275$ και $275\cdot1{,}1=302{,}5$ €. Συνολική αύξηση $\frac{52{,}5}{250}=0{,}21=21\%$ (όχι $20\%$).
\sol{10} Έτος $1$: τόκος $100$ €, κεφάλαιο $5\,100$. Έτος $2$: τόκος $102$ €. Σύνολο τόκων $202$ € και ποσό $5\,202$ €.
\sol{11} $0{,}03x=60$, $x=2\,000$ €.
\sol{12} $12+0{,}05\cdot800=52$ €, με ΦΠΑ $1{,}24\cdot52=64{,}48$ €.
\sol{13} Μετρητά $200$ €, υπόλοιπο $600$ €. Τόκος $3\cdot0{,}01\cdot600=18$ €, δηλαδή $\frac{18}{800}=2{,}25\%$ της αρχικής τιμής.
\sol{14} Ύπνος $33{,}3\%$, σχολείο $25\%$, διάβασμα $12{,}5\%$, αθλητισμός $8{,}3\%$, υπόλοιπα $20{,}8\%$. Μήκη ράβδων: $3{,}3$, $2{,}5$, $1{,}25$, $0{,}83$ και $2{,}08$ εκ. (το άθροισμα βγαίνει $10$ εκ. με στρογγυλοποίηση).
\sol{15} Η νέα τιμή είναι $0{,}8$ της παλιάς, και $1:0{,}8=1{,}25$. Πρέπει να ανέβει κατά $25\%$.
\sol{16} $100\to70\to56$. Η τιμή έχει μειωθεί κατά $44$ €, δηλαδή $44\%$.
\sol{17} Όχι: $100\to50\to75$, δηλαδή $25\%$ χαμηλότερα.
\sol{18} Πληρώνεις $2$ από τα $3$, άρα η έκπτωση είναι $\frac{1}{3}\approx33{,}3\%$.
\sol{19} $100-28-35-12=25\%$.
\sol{20} Με αξία $100$ € ο ΦΠΑ είναι $24$ € και η τιμή $124$ €. Το $24$ είναι τα $\frac{24}{124}\approx19{,}4\%$ της τελικής τιμής.
\sol{21} Ελεύθερη απάντηση.
\sol{22} Ελεύθερη απάντηση. Αθροίζουν τα ποσοστά σε $100\%$.
